\documentclass{beamer}
\usepackage{amsmath, tikz, mathabx, MnSymbol}
\usetikzlibrary{arrows.meta,graphs,graphdrawing}
\usegdlibrary{layered}
\nonstopmode

\usetheme{Antibes}
\usecolortheme{beaver}

\title{Optimisation des gains au Blackjack}
\author{Victor Lajoux-Thiebault}
\date{2026}

\begin{document}

\section{Introduction}

\begin{frame}
	\titlepage
\end{frame}

\begin{frame}{Outline}
	\setcounter{tocdepth}{2}
	\tableofcontents
\end{frame}

\section{Étude d'une manche}
\subsection{Définitions}
\subsubsection{Variables du joueur}

\begin{frame}{Variables du joueur}
	On utilisera les notations suivantes:
	\begin{itemize}
		\item $p \in \mathbb{R}^{10} $ est le vecteur associé à la pioche. Chaque i-ème coordonnée représente la quantité de cartes avec une valeur i dans la pioche. \\{} e.g.  
		La pioche à l'origine est $d = \left(\begin{smallmatrix} 4 \\ 4 \\ 4 \\ 4 \\ 4 \\ 4 \\ 4 \\ 4 \\ 4 \\ 16 \end{smallmatrix}\right)$.
		\item $j_1, j_2 \in \mathbb{R}^{10}$ de la même manière sont les vecteur représentant la première main du joueur, et dans le cas d'un split, sa deuxième main.
		\item $c \in \mathbb{R}^{10}$ représente la main du croupier. \\{} 
		Lorsque $|c| := \sum c_i = 1$, il est noté $c^{*}$.
	\end{itemize}
\end{frame}

\subsubsection{États du jeu}

\begin{frame}{États du jeu}
	\begin{itemize}
		\item $A = \{\text{rester},\text{tirer}, \text{doubler}, \text{split} \}$ est l'ensemble des actions faisable par le joueur. Lorsque les règles forcent certaines actions, nous admetterons que le joueur joue l'action associé (souvent: Rester).
		\item Une situation/état $s$ est un 4-uplet $(p, j_1, j_2, c, d, m)$ où $d$ est un 2-uplet de booléen indiquant si le joueur a doublé sur sa i-ème main ou pas, et $m \in \{1,2\}$ indiquant quelle main est en train d'être joué dans l'éventualité d'un split. Les états sont relié par les actions de $A$.
		\item $S$ est l'ensemble des états du jeu.
	\end{itemize}
\end{frame}

\begin{frame}
	On utilisera aussi les notations suivantes:
	\begin{itemize}
		\item $S^{p_0}$ l'ensemble des situations atteignables à partir d'une pioche d'origine $p_0$
		\item $S(s_0)$ l'ensemble des situations atteignable à partir de l'état $s_0$.
		\item $S_t^{p_0}$ l'ensemble des situations atteignable en effectuant $t$ actions à partir de la pioche d'origine $p_0$.
		\item $S^F$ l'ensemble des situations où $c$ a au moins 16 points.
	\end{itemize}
\end{frame}

\subsubsection{Fonction de transition probabiliste}

\begin{frame}{Fonction de transition probabiliste}

	\begin{itemize}
		\item Pour tout $(s,s',a) \in S\times S\times A$ on note $T_a(s,s')$ la probabilité de passer de la situation $s$ à $s'$ en effectuant l'action $a$.
		\item Si $a \neq$ Split, $$T_a(s,s') = \frac{|p\cdot (p - p')|}{|p|}$$
			sinon, $$T_{Split}(s,s') = 1$$ car il n'y a qu'une seule situation atteignable.
	\end{itemize}
\end{frame}


\begin{frame}
	Illustration à compléter.
	\tikz [>={Stealth[round,sep]}]
	\graph [layered layout, components go right top aligned, nodes={draw, circle}, edges=rounded corners]
		{
			s1 -> {rester[fill=red!50, draw=red, radius=5] -> {s2, s3, s4}, tirer[fill=red!50, draw=red, radius=5] -> {s5, s6, s7}, doubler[fill=red!50, draw=red, radius=5] -> {s5', s6', s7'}};
		};
\end{frame}

\subsubsection{Manche de jeu}

\begin{frame}{Manche de jeu}
	\begin{itemize}
		\item Une stratégie est une fonction $\theta : S \rightarrow A$ qui définit l'action du joueur pour chaque situation.
		\item Une manche $\mathcal{M}_{\theta}^{p_0}$ est une suite $(s_t)_{0\leq t \leq n}$ tel que $s_t \in S_t^{p_0}$ et $s_n$ est un état final.
	\end{itemize}
\end{frame}

\subsubsection{Variable aléatoire d'intérêt}

\begin{frame}{Variable aléatoire d'intérêt}
	Pour toute situation $s \in S^{p_0}$ on s'intéresse à la variable aléatoire qui renvoie le gain d'une manche $\mathcal{M}^{p_0}_\theta$ commençant par $s$. On la note : $$
	Y_\theta^{(s=s_t)} = \begin{cases}
		\mathcal{R}(s_t) \text{ si } s_t \in S^F \\
		Y_\theta^{s_{t+1}} \text{ sinon}
	\end{cases}
	$$
	où $\mathcal{R}(s_t) \in \{-2, -1, 0, 1, 1.5, 2\}$ selon les règles établient plus tôt. \\
	
	En particulier, pour un état initial $s_0$, celui-ci est entièrement déterminé par $p_0$, on note ainsi:
	$$ X_\theta^{p_0} = Y_\theta^{s_0} $$

\end{frame}

\subsection{Optimisation}

\subsubsection{Mise en équation du problème}

\begin{frame}{Mise en équation du problème}
	Le but du problème est alors pour une pioche $p_0$ donnée, trouver $\theta^*$ tel que $\mathbb{E}[X_{\theta^*}^{p_0}]$ est maximal.
	L'existence d'une stratégie maximale est donnée par le caractére fini de l'ensemble $S \times A$.
\end{frame}

\begin{frame}
	Soit $\theta, p_0$ et $s_t$. 
	\[
		\mathbb{E}[Y_\theta^{s_t}] &= \sum_{\substack{(s_t, \dots, s_n) \\ \text{une manche}}} \mathbb{P}(s_{t+1}, \dots, s_n \,|\, s_t) \mathcal{R}(s_n)
	\]
	Or, par distributivité de l'union sur le produit cartésien
	\begin{align*}
		\bigtimes_{k=t+1}^n S_k &= \left( \bigcupdot_{s \in S_{t+1}} \{s\} \right) \bigtimes_{k=t+2}^n S_k
		&= \bigcupdot_{s \in S_{t+1}} \left( \{s\} \bigtimes_{k=t+2}^n S_k \right)
	\end{align*}
	donc par la formule des probabilités totales
	\[ 
		\mathbb{P}(s_{t+1}, \dots, s_n) = \sum_{s \in S_{t+1}} \mathbb{P}(s_{t+1} = s) \mathbb{P}(s_{t+2}, \dots, s_n \,|\,  s_t, s_{t+1} = s)
	\]

\end{frame}

\begin{frame}
	Ainsi, les sommes étant bien finies
	\begin{align*}
		\mathbb{E}[Y_\theta^{s_t}] &= \sum_{\substack{(s_t, \dots, s_n) \\ \text{une manche}}} \sum_{s \in S_{t+1}} \mathbb{P}(s_{t+1} = s) \mathbb{P}(s_{t+2}, \dots, s_n \,|\, s_{t+1}) \mathcal{R}(s_n) \\
		&=  \sum_{s \in S_{t+1}} \mathbb{P}(s \,|\, s_t) \sum_{\substack{(s_t, s, s_{t+2} \dots, s_n) \\ \text{une manche}}} \mathbb{P}(s_{t+2}, \dots, s_n \,|\, s_{t+1}) \mathcal{R}(s_n) \\
		&= \sum_{s \in S_{t+1}} \mathbb{P}(s \,|\, s_t) \mathbb{E}[Y_\theta^{s_{t+1}}] \\
		&= \sum_{s \in S_{t+1}} T_{\theta(s_t)}(s_t, s) \mathbb{E}[Y_\theta^{s_{t+1}}]
	\end{align*}
\end{frame}

\begin{frame}
	D'où la relation de récurrence suivante : 
	\[
		\mathbb{E}[Y_\theta^{s_t}] &= \sum_{s \in S_{t+1}} T_{\theta(s_t)}(s_t, s) \mathbb{E}[Y_\theta^{s_{t+1}}]
	\]
	On remarque alors que $\mathbb{E}[Y_{\theta}^{s_t}]$ est maximale seulement si tout les $\mathbb{E}[Y_{\theta}^{s_{t+1}}]$ sont maximales.
\end{frame}

\begin{frame}
	On en déduit les relations de récurrences : 
		\[
			\boxed{
			\theta^*(s_t) = \text{arg}\max_{a \in A} \sum_{s \in S_{t+1}} T_{a}(s_t, s) \mathbb{E}[Y_{\theta^*}^{s_{t+1}}]
			}
		\]
		et 
		\[	
			\boxed{
			\mathbb{E}[Y_{\theta^*}^{s_t}] &= \max_{a \in A} \sum_{s \in S_{t+1}} T_{a}(s_t, s) \mathbb{E}[Y_{\theta^*}^{s_{t+1}}]
			}
		\]
			
\end{frame}

\subsection{Implémentation informatique}

\subsubsection{Analyse de la complexité}

\begin{frame}{Analyse de la complexité}
	Il suffit alors de calculer $\theta^*$ de manière récursive. Mais alors, combiens il y a t-il d'appel récursif?
\end{frame}

\begin{frame}
	Pour $s \in S$, quelle est le cardinal de $ S_a(s) $ ?
	On suppose la présence d'au moins une carte de chaque valeur dans la pioche.
\end{frame}

\end{document}
